% Introduction chapter

\section{Background}
The Baltic Sea is a critical maritime region with significant commercial and strategic importance. However, the presence of shadow fleet vessels—ships that operate outside normal regulatory frameworks and tracking systems—poses challenges for maritime security, environmental protection, and international trade monitoring.

\section{Problem Statement}
Shadow fleet vessels often exhibit distinctive behavioral patterns that differentiate them from legitimate commercial traffic. These patterns include irregular routes, AIS (Automatic Identification System) signal manipulation, frequent flag changes, and operations in areas known for sanctions evasion or illegal activities.

Traditional manual monitoring methods are time-consuming and cannot scale to the volume of vessel traffic in the Baltic Sea. This creates a need for automated detection systems capable of identifying suspicious vessel behavior.

\section{Research Questions}
This thesis addresses the following research questions:
\begin{enumerate}
    \item How can machine learning techniques be applied to detect shadow fleet vessels based on their movement patterns and behavioral characteristics?
    \item What features from vessel tracking data are most indicative of shadow fleet operations?
    \item What level of accuracy can be achieved in identifying shadow fleet vessels using the proposed approach?
\end{enumerate}

\section{Objectives}
The main objectives of this thesis are:
\begin{itemize}
    \item Develop a machine learning model for shadow fleet vessel detection
    \item Identify and extract relevant features from vessel tracking data
    \item Evaluate the model's performance using appropriate metrics
    \item Provide insights into the behavioral patterns of shadow fleet vessels
\end{itemize}

\section{Scope and Limitations}
This work focuses on:
\begin{itemize}
    \item Vessels operating in the Baltic Sea region
    \item AIS data and publicly available vessel information
    \item Pattern recognition and classification techniques
\end{itemize}

Limitations include:
\begin{itemize}
    \item Availability and quality of training data on known shadow fleet vessels
    \item AIS data gaps and potential spoofing
    \item Evolving tactics used by shadow fleet operators
\end{itemize}

\section{Thesis Outline}
The remainder of this thesis is organized as follows:
\begin{itemize}
    \item Chapter~\ref{ch:background} provides background on maritime tracking systems, shadow fleet operations, and related work in vessel behavior analysis.
    \item Chapter~\ref{ch:methodology} describes the methodology, including data collection, feature engineering, and machine learning algorithms used.
    \item Chapter~\ref{ch:implementation} details the implementation of the detection system.
    \item Chapter~\ref{ch:results} presents the results of the experiments and model evaluation.
    \item Chapter~\ref{ch:discussion} discusses the findings, their implications, and limitations.
    \item Chapter~\ref{ch:conclusions} concludes the thesis and suggests directions for future work.
\end{itemize}
